% 3-3-1-benchmark.tex
%  Budget: 3pp.
%  Source map: protocol -> lago_forecasting_2021 (open benchmark, years_test=2).
%  ALL numbers computed from data/raw/DE.csv via BenchmarkLoader on 2026-09-05
%  and cross-checked against HANDOFF §19.0.2:
%    full  2012-01-09..2017-12-31, n=52416
%    train 2012-01-09..2016-01-03, n=34944, mean 36.1976, std 15.9463
%    test  2016-01-04..2017-12-31, n=17472, 728 origins, mean 31.6383
%    negatives: 538 full / 297 train / 241 test; min -221.99, max 210.00
%  Rules: \lr{} every numeral; \lr{—} for em-dash.

\subsection{مجموعه‌داده بنچمارک \lr{EPEX-DE}}

مقایسه اصلی این پژوهش بر مجموعه‌داده باز بازار آلمان انجام شده است؛ همان
مجموعه‌ای که لاگو و همکاران \cite{lago_forecasting_2021} به‌عنوان بخشی از
بنچمارک پیشنهادی خود منتشر کردند. انتخاب این مجموعه‌داده، تصمیم
روش‌شناختی محوری فصل حاضر است و دو دلیل دارد.

نخست، عمومی بودن. مجموعه‌داده و کد مرجع هر دو در دسترس‌اند، بنابراین
اعداد این پایان‌نامه را می‌توان بازتولید کرد. این دقیقاً همان چیزی است که
نقد بخش \lr{2-9} \lr{—} اتکای ادبیات به داده‌های منتشرنشده \lr{—} به آن
اشاره داشت.

دوم، وجود اعداد منتشرشده بر همان داده. این امکان مقایسه مستقیم را فراهم
می‌کند، و بخش \lr{4-5-2} از آن استفاده می‌کند. بدون این ویژگی، هر ادعای
«بهتر بودن» بی‌معنا می‌شد.

\subsubsection*{ساختار و دامنه}

مجموعه‌داده سه سری ساعتی دارد و نه بیشتر:

\begin{itemize}
  \item \lr{price} \lr{—} قیمت تسویه روز-پیش بازار \lr{EPEX-DE}، بر حسب
        یورو بر مگاوات‌ساعت.
  \item \lr{exog\_1} \lr{—} پیش‌بینی روز-پیش بار منتشرشده توسط
        \lr{Amprion}.
  \item \lr{exog\_2} \lr{—} پیش‌بینی روز-پیش مجموع تولید بادی و خورشیدی.
\end{itemize}

هر دو متغیر برون‌زا به شکل \emph{پیش‌بینی} و نه مقدار تحقق‌یافته منتشر
می‌شوند. این تمایز، پایه مجاز بودن استفاده از آن‌ها در لحظه مرجع است و در
بخش \lr{3-4} به آن بازمی‌گردیم.

دامنه زمانی کامل از \lr{2012-01-09} تا \lr{2017-12-31} است و شامل
\lr{52416} ساعت می‌شود.

\subsubsection*{تقسیم آموزش و آزمون}

مطابق پروتکل مرجع، دو سال پایانی به‌عنوان دوره آزمون کنار گذاشته می‌شود.
این تقسیم، پیش از هر مدل‌سازی و صرفاً بر پایه تقویم انجام شده است و در
هیچ مرحله‌ای بازبینی نشده است:

\begin{center}
\begin{latin}
\begin{tabular}{lccc}
\hline
 & span & hours & origins \\
\hline
train & 2012-01-09 \dots 2016-01-03 & 34,944 & --- \\
test  & 2016-01-04 \dots 2017-12-31 & 17,472 & 728 \\
\hline
\end{tabular}
\end{latin}
\end{center}

دوره آزمون \lr{728} مبدأ پیش‌بینی روزانه در بر می‌گیرد. پنجره
اعتبارسنجی، که برای تنظیم ابرپارامترها به کار می‌رود، از انتهای دوره
آموزش جدا می‌شود و اکیداً پیش از آغاز دوره آزمون پایان می‌یابد
(بخش \lr{3-5}).

\subsubsection*{آمار توصیفی قیمت}

\begin{center}
\begin{latin}
\begin{tabular}{lrrrrr}
\hline
 & mean & std & min & max & negative hours \\
\hline
full  & 34.6779 & 15.9400 & -221.99 & 210.00 & 538 \\
train & 36.1976 & 15.9463 & -221.99 & 210.00 & 297 \\
test  & 31.6383 & 15.4865 & -130.09 & 163.52 & 241 \\
\hline
\end{tabular}
\end{latin}
\end{center}

سه نکته از این جدول برای فصول بعد اهمیت دارد.

\textbf{قیمت‌های منفی واقعی‌اند و کم‌شمار نیستند.} در کل دوره،
\lr{538} ساعت با قیمت منفی ثبت شده است، یعنی حدود \lr{1.03} درصد ساعات، و
کمینه مطلق به \lr{-221.99} یورو بر مگاوات‌ساعت می‌رسد. این مشاهده به‌
تنهایی معیارهای خطای نسبی مبتنی بر تقسیم بر مقدار واقعی را کنار می‌گذارد،
و دلیل کنارگذاشتن \lr{MAPE} در بخش \lr{3-5} است.

\textbf{میانگین دوره آزمون از دوره آموزش پایین‌تر است}
(\lr{31.6383} در برابر \lr{36.1976}). بنابراین حتی درون همین مجموعه‌داده،
مدل بر دوره‌ای با سطح قیمت متفاوت آزموده می‌شود. این تفاوت کوچک است، اما
هم‌جهت با پدیده‌ای است که در آزمون برون‌توزیع فصل پنجم به‌مراتب
بزرگ‌تر ظاهر می‌شود.

\textbf{انحراف معیار دوره آموزش} (\lr{15.9463}) مبنای تعریف آستانه رژیم
در بخش \lr{3-8} است. آستانه از میانگین و انحراف معیار داده‌ای پیش از
پنجره‌های تنظیم و وزن‌دهی محاسبه می‌شود، نه از کل دوره آموزش؛ جزئیات و
دلیل این سخت‌گیری در همان بخش می‌آید.

\subsubsection*{بازتولیدپذیری}

داده خام یک بار دریافت و به‌صورت یک تصویر ثابت با اثر انگشت
\lr{sha256} در مخزن ثبت شده است، و کد ویژگی‌سازی تنها از همان تصویر
می‌خواند. دلیل این سخت‌گیری ساده است: یک واسط برنامه‌نویسی زنده، فردا
پاسخ دیگری می‌دهد و این کار را بی‌صدا انجام می‌دهد. جداسازی شبکه از
مسیر ساخت ویژگی‌ها با آزمون خودکار تضمین می‌شود، نه با قرارداد.
